%%
%% fall-lecture03.tex
%% 
%% Made by Alex Nelson <pqnelson@gmail.com>
%% Login   <alex@lisp>
%% 
%% Started on  2026-10-08T09:29:44-0700
%% Last update 2026-10-08T09:29:44-0700
%% 

\lecture{}

\begin{node}
The problem set will be Folland, Chapter 4, problems: 6, 14, 20, 26, 32.
Today we will cover roughly Folland~\S4.3.
\end{node}

\begin{recall}
We saw how sequences were inadequate to fully describe topologies. But
there is hope for ``generalized sequences'' (nets) could do this.
\end{recall}

\begin{definition}
A \define{Directed Set} $(\Lambda,\lesssim)$ is a paritally ordered
set [i.e., transitive $\lambda_{1}\lesssim\lambda_{2}$ and
  $\lambda_{2}\lesssim\lambda_{3}$ implies $\lambda_{1}\lesssim\lambda_{3}$,
  and reflexive $\lambda\lesssim\lambda$] satisfying:
\begin{enumerate}
\item for every $\lambda_{1}$, $\lambda_{2}\in\Lambda$, there exists a
  $\lambda\in\Lambda$ such that $\lambda_{1}\lesssim\lambda$ and
  $\lambda_{2}\lesssim\lambda$ [i.e., $\lambda$ is even bigger than
  both $\lambda_{1}$ and $\lambda_{2}$]. 
\end{enumerate}
\end{definition}

\begin{example}
We see $\Lambda=\NN$ with the usual ordering is a directed set.
\end{example}

\begin{example}
Let $(A,\lesssim)$ and $(B,\lesssim)$ be two directed sets.
We see that $\Lambda=A\times B$ with $(a,b)\lesssim(a',b')$ iff
$a\lesssim a'$ and $b\lesssim b'$. This is a directed set.
\end{example}

\begin{example}
Let $(X,\topology)$ be a topological space. Let $x\in X$ be a point.
Define $\Lambda$ as some ``neighborhood base'' at $x$ (so $\Lambda$ is
a collection of neighborhoods of $x$ such that for every neighborhood
$U$ of $x$, there exists a $V\in\Lambda$ such that $V\subset U$). Then
$\lambda_{1}\lesssim\lambda_{2}$ iff $U_{\lambda_{2}}\subset U_{\lambda_{1}}$
(this is directed by reverse inclusion).
\end{example}

\begin{definition}
Let $(X,\topology)$ be a topological space. A \define{Net} is a map
$\Lambda\to X$ where $\Lambda$ is a directed set. The notation for a
net is usually written as $\{x_{\lambda}\}_{\lambda\in\Lambda}$.
\end{definition}

\begin{definition}
Let $(X,\topology)$ be a topological space.
Let $\{x_{\lambda}\}_{\lambda\in\Lambda}$ be a net on $X$. We say
$x_{\lambda}$ \define{Converges to} $x$ (and write it as
``$x_{\lambda}\to x$'') if for every neighborhood $U\ni x$ of $x$,
there exists a $\lambda_{0}\in\Lambda$ such that for all
$\lambda\in\Lambda$ we have $\lambda\geq\lambda_{0}$ implies
$x_{\lambda}\in U$.
\end{definition}

\begin{remark}
Observe that a net defined on the directed set $\NN$ is precisely the
usual notion of a sequence we all know and love. Moreover, we recover
the usual notion of convergence of sequences when we look at
``convergence of nets''.
\end{remark}

\begin{theorem}
For any topological space $(X,\topology)$, for any subset $E\subset X$,
for any point $x\in X$, we have $x\in\closure{E}$ if and only
if there exists a net $\{x_{\lambda}\}_{\lambda\in\Lambda}$ in $E$
such that $x_{\lambda}\to x$ (the net converges to $x$).
\end{theorem}

\begin{theorem}
If $f\colon X\to Y$ is a function between topological spaces,
then $f$ is continuous at $x\in X$ if and only if $x_{\lambda}\to x$
implies $f(x_{\lambda})\to f(x)$.
\end{theorem}

\begin{theorem}
We have $f_{\lambda}\to f$ is a convergent net in $\Maps(X,Y)=Y^{X}$ the
space of maps if and only if for all $x\in X$ we have
$f_{\lambda}(x)\to f(x)$ converges pointwise.
\end{theorem}

\subsection{Countability axioms of topological spaces}

\begin{definition}
A topological space $(X,\topology)$ is \define{First Countable}
if every point has a countable neighborhood base. We indicate this by
writing $(X,\topology)\in(A1)$.
\end{definition}

\begin{example}
Every metric space is first-countable. Just consider the open balls
$B(x,1/n)$ for $n\in\NN$ as the base.
\end{example}

\begin{remark}
First-countability is a \emph{local} property, i.e., considered about
a point (for every point).
\end{remark}

\begin{theorem}
Let $(X,\topology)$ be a \emph{first-countable} space. Let $E\subset X$.
Let $x\in X$. Then $x\in\closure{E}$ if and only if there exists a
(``standard'') sequence $(x_{n})\subset E$ such that $x_{n}\to x$.
\end{theorem}

\begin{definition}
We call the topological space $(X,\topology)$ \define{Second Countable}
if there exists a countable base for the topology. We sometimes write
this as $(X,\topology)\in(A2)$.
\end{definition}

\begin{theorem}
Second-countable spaces are first-countable.
\end{theorem}

\begin{lemma}
A subset $\beta\subset\topology$ (of a topology $\topology$) is a base
if and only if $\beta$ contains a neighborhood base at every point.
\end{lemma}

\begin{example}
Let $X=\ell^{\infty}(\NN)$ consist of bounded sequences
$(a_{1},a_{2},\dots)$. Let us use the sup-metric
\begin{equation}
\rho(a,b)=\sup_{n}|a_{n}-b_{n}|.
\end{equation}
Then $(X,\rho)$ is a metric space, so it is first-countable. But it is
\emph{not} second-countable!
\end{example}

\begin{proof}
We want to show $\ell^{\infty}$ does not have a countable base for its
metric topology. Consider, for every $A\subset\NN$, the indicator
function on $A$. The set $2^{\NN}$ is uncountable (duh). Then the set
of $(x_{n})$ where $x_{n}=\mathbf{1}_{A}(n)$ for each $A\subset\NN$,
and the distance between any two such $A,B\subset\NN$ where $A\neq B$
is $1$. Then we have an uncountable set of elements of $\ell^{\infty}$
which are a distance 1 far from each other. This cannot happen for
second-countable spaces: if $\ell^{\infty}$ were second-countable,
there would be a countable dense set, rom which you can find distinct
arbitrarily close points from the dense set $\{\mathbf{1}_{A}\mid A\subset\NN\}$,
which gives an injective function from an uncountable set to a
countable set, which is a contradiction.
\end{proof}
